\lecture{09}{Условие Слейтера и сильная двойственность}

\sourcevideo
    {https://www.youtube.com/playlist?list=PLOzRYVm0a65fhKDS8cYIC7YCuVGJ4FOJx}
    {Week 3: Lecture 9: Slater's Condition}

Для прямой задачи с оптимальным значением \(p^\star\) и её
лагранжевой двойственной задачи с оптимальным значением \(d^\star\)
слабая двойственность даёт
\[
    d^\star\le p^\star.
\]
Разность между этими значениями может быть положительной. Условие
Слейтера является основным достаточным условием, при котором
двойственный разрыв исчезает.

\section{Двойственный разрыв}

Рассмотрим прямую задачу
\[
    \begin{aligned}
        \minimize_x
        \quad&
        f_0(x),
        \\
        \subjectto
        \quad&
        f_i(x)\le0,
        \qquad i=1,\ldots,m,
        \\
        &
        h_j(x)=0,
        \qquad j=1,\ldots,r.
    \end{aligned}
\]
Её функция Лагранжа и двойственная функция равны
\[
    \mathcal L(x,\lambda,\nu)
    =
    f_0(x)
    +
    \sum_{i=1}^{m}\lambda_i f_i(x)
    +
    \sum_{j=1}^{r}\nu_jh_j(x),
\]
\[
    g(\lambda,\nu)
    =
    \inf_x\mathcal L(x,\lambda,\nu).
\]

Двойственная задача имеет вид
\[
    d^\star
    =
    \sup_{\lambda\ge0,\nu}g(\lambda,\nu).
\]

\begin{definition}[Двойственный разрыв]
Если оптимальные значения \(p^\star\) и \(d^\star\) конечны, то
величина
\[
    \Delta_{\mathrm{dual}}
    =
    p^\star-d^\star
    \ge0
\]
называется двойственным разрывом.
\end{definition}

Если
\[
    p^\star=d^\star,
\]
выполняется сильная двойственность. В противном случае двойственная
задача даёт лишь строгую нижнюю оценку прямого оптимального значения.

\section{Пример положительного двойственного разрыва}

Рассмотрим задачу
\[
    \begin{aligned}
        \minimize_{x=(x_1,x_2)\in\RR^2}
        \quad&
        \exp(x_2),
        \\
        \subjectto
        \quad&
        \norm{x}_2\le x_1.
    \end{aligned}
    \tag{\ensuremath{P}}
\]

Ограничение эквивалентно
\[
    \sqrt{x_1^2+x_2^2}\le x_1.
\]
Обе части неотрицательны, поэтому после возведения в квадрат
получаем
\[
    x_1^2+x_2^2\le x_1^2,
\]
откуда
\[
    x_2=0,
    \qquad
    x_1\ge0.
\]
Следовательно, на всём допустимом множестве
\[
    \exp(x_2)=1,
\]
и потому
\[
    p^\star=1.
\]

Запишем ограничение как
\[
    h(x)
    =
    \norm{x}_2-x_1
    \le0.
\]
Функция Лагранжа имеет вид
\[
    \mathcal L(x,\lambda)
    =
    \exp(x_2)
    +
    \lambda\bigl(\norm{x}_2-x_1\bigr),
    \qquad
    \lambda\ge0,
\]
а двойственная функция равна
\[
    g(\lambda)
    =
    \inf_{x\in\RR^2}
    \left[
        \exp(x_2)
        +
        \lambda\bigl(\norm{x}_2-x_1\bigr)
    \right].
\]

Поскольку
\[
    \exp(x_2)>0,
    \qquad
    \norm{x}_2-x_1\ge0
\]
для любого \(x\), имеем
\[
    g(\lambda)\ge0.
\]
Следовательно,
\[
    d^\star\ge0.
\]

Теперь ограничим поиск последовательностью
\[
    x_1=x_2^4,
    \qquad
    x_2\to-\infty.
\]
Тогда
\[
\begin{aligned}
    \norm{x}_2-x_1
    &=
    \sqrt{x_2^8+x_2^2}-x_2^4
    \\
    &=
    \frac{x_2^2}{
        \sqrt{x_2^8+x_2^2}+x_2^4
    }
    \longrightarrow0.
\end{aligned}
\]
Одновременно
\[
    \exp(x_2)\longrightarrow0.
\]
Поэтому для каждого \(\lambda\ge0\)
\[
    g(\lambda)\le0.
\]
Вместе с ранее полученной оценкой это даёт
\[
    g(\lambda)=0,
    \qquad
    d^\star=0.
\]

Таким образом,
\[
    p^\star=1,
    \qquad
    d^\star=0,
    \qquad
    p^\star-d^\star=1.
\]
В этой задаче сильная двойственность не выполняется.

\section{Строгая допустимость}

Пусть прямая задача выпукла: функция \(f_0\) и ограничения
\(f_i\) выпуклы, а равенства имеют аффинный вид
\[
    h_j(x)=a_j^{\mathsf T}x-b_j.
\]

\begin{definition}[Строго допустимая точка]
Точка \(\bar x\) называется строго допустимой, если
\[
    f_i(\bar x)<0,
    \qquad i=1,\ldots,m,
\]
и одновременно
\[
    h_j(\bar x)=0,
    \qquad j=1,\ldots,r.
\]
\end{definition}

Строгость требуется только для неравенств:
\[
    f_i(\bar x)<0.
\]
Аффинные равенства должны выполняться точно.

\section{Условие Слейтера}

\begin{theorem}[Условие Слейтера]
Рассмотрим выпуклую задачу
\[
    \begin{aligned}
        \minimize_x
        \quad&
        f_0(x),
        \\
        \subjectto
        \quad&
        f_i(x)\le0,
        \qquad i=1,\ldots,m,
        \\
        &
        Ax=b.
    \end{aligned}
\]
Пусть \(p^\star\) конечно и существует точка
\[
    \bar x\in
    \operatorname{ri}
    \bigl(
        \operatorname{dom}f_0
        \cap
        \textstyle\bigcap_i\operatorname{dom}f_i
    \bigr),
\]
для которой
\[
    f_i(\bar x)<0,
    \qquad
    A\bar x=b.
\]
Тогда выполняется сильная двойственность:
\[
    p^\star=d^\star.
\]
При этих предположениях существует и оптимальная двойственная точка
\((\lambda^\star,\nu^\star)\).
\end{theorem}

Здесь \(\operatorname{ri}\) обозначает относительную внутренность.
Она используется потому, что допустимое множество может лежать в
аффинном подпространстве меньшей размерности.

В рассмотренном выше примере допустимое множество имеет вид
\[
    \setgiven{(x_1,x_2)}{x_1\ge0,\ x_2=0}.
\]
Для каждой его точки
\[
    \norm{x}_2-x_1=0,
\]
поэтому строгого неравенства получить нельзя. Условие Слейтера не
выполняется, что согласуется с положительным двойственным разрывом.

\begin{remark}[Логика условия]
Условие Слейтера является достаточным, но не необходимым. Из его
невыполнения нельзя заключить, что сильная двойственность обязательно
нарушается. Однако если сильная двойственность нарушена, условия
теоремы Слейтера выполнены быть не могут.
\end{remark}

\begin{remark}[Аффинные ограничения]
В уточнённой форме условия Слейтера строгие неравенства требуются
только для неаффинных выпуклых ограничений. Аффинные неравенства
достаточно выполнять нестрого. Поэтому, если все ограничения аффинны,
задача допустима и её оптимальное значение конечно, сильная
двойственность выполняется в стандартной конечномерной постановке.
\end{remark}

\section{Следствия сильной двойственности}

Если
\[
    p^\star=d^\star,
\]
прямую задачу можно решать через двойственную, не теряя оптимального
значения. Это особенно полезно, когда размерность двойственной
переменной меньше размерности прямой переменной.

Кроме того, при выполнении подходящего условия регулярности условия
Каруша--Куна--Таккера становятся не только достаточными, но и
необходимыми для оптимальности выпуклой задачи. Эти условия будут
введены отдельно.

\section{Чувствительность оптимального значения}

Рассмотрим параметрическую задачу
\[
    p(b)
    =
    \minimize_x
    \left\{
        f(x)
        \,\middle|\,
        Bx=b
    \right\}.
\]
Её функция Лагранжа равна
\[
    \mathcal L(x,\nu;b)
    =
    f(x)+\nu^{\mathsf T}(Bx-b).
\]

Пусть решение \(x^\star(b)\) и оптимальный множитель
\(\nu^\star(b)\) определены однозначно и достаточно гладко зависят
от \(b\). Тогда
\[
    p(b)
    =
    \mathcal L
    \bigl(
        x^\star(b),
        \nu^\star(b);
        b
    \bigr).
\]

Дифференцируя и используя условия стационарности
\[
    \nabla_x\mathcal L=0,
    \qquad
    \nabla_\nu\mathcal L=0,
\]
получаем
\[
    \nabla p(b)
    =
    \nabla_b\mathcal L
    =
    -\nu^\star(b).
\]

\begin{remark}[Знак множителя]
При соглашении
\[
    \mathcal L(x,\nu;b)
    =
    f(x)+\nu^{\mathsf T}(Bx-b)
\]
справедливо
\[
    \nabla p(b)=-\nu^\star.
\]
При записи ограничения в форме \(b-Bx=0\) знак меняется. Поэтому
экономическая интерпретация множителя зависит от выбранного
соглашения о знаках.
\end{remark}

Двойственная переменная характеризует чувствительность оптимального
значения к изменению правой части ограничения. Поэтому множители
Лагранжа также называют теневыми ценами или предельными стоимостями.

\section{Двойственная форма квадратичной задачи}

Рассмотрим задачу
\[
    \begin{aligned}
        \minimize_{x\in\RR^n}
        \quad&
        \frac12x^{\mathsf T}Qx+c^{\mathsf T}x,
        \\
        \subjectto
        \quad&
        Ax\le b,
    \end{aligned}
    \tag{\ensuremath{P_Q}}
\]
где
\[
    Q=Q^{\mathsf T}\succ0,
    \qquad
    A\in\RR^{m\times n}.
\]

Функция Лагранжа равна
\[
\begin{aligned}
    \mathcal L(x,\lambda)
    &=
    \frac12x^{\mathsf T}Qx
    +
    c^{\mathsf T}x
    +
    \lambda^{\mathsf T}(Ax-b)
    \\
    &=
    \frac12x^{\mathsf T}Qx
    +
    (c+A^{\mathsf T}\lambda)^{\mathsf T}x
    -
    b^{\mathsf T}\lambda,
\end{aligned}
\]
где
\[
    \lambda\ge0.
\]

Для вычисления двойственной функции минимизируем
\(\mathcal L\) по \(x\). Условие стационарности имеет вид
\[
    Qx+c+A^{\mathsf T}\lambda=0.
\]
Поскольку \(Q\succ0\),
\[
    x(\lambda)
    =
    -Q^{-1}
    \bigl(
        c+A^{\mathsf T}\lambda
    \bigr).
\]

Подставляя это выражение в функцию Лагранжа, получаем
\[
    g(\lambda)
    =
    -
    \frac12
    \bigl(
        c+A^{\mathsf T}\lambda
    \bigr)^{\mathsf T}
    Q^{-1}
    \bigl(
        c+A^{\mathsf T}\lambda
    \bigr)
    -
    b^{\mathsf T}\lambda.
\]

Следовательно, двойственная задача равна
\[
    \begin{aligned}
        \maximize_{\lambda\in\RR^m}
        \quad&
        -
        \frac12
        \bigl(
            c+A^{\mathsf T}\lambda
        \bigr)^{\mathsf T}
        Q^{-1}
        \bigl(
            c+A^{\mathsf T}\lambda
        \bigr)
        -
        b^{\mathsf T}\lambda,
        \\
        \subjectto
        \quad&
        \lambda\ge0.
    \end{aligned}
    \tag{\ensuremath{D_Q}}
\]

После раскрытия скобок
\[
\begin{aligned}
    g(\lambda)
    ={}&
    -\frac12
    \lambda^{\mathsf T}
    AQ^{-1}A^{\mathsf T}
    \lambda
    \\
    &-
    \bigl(
        AQ^{-1}c+b
    \bigr)^{\mathsf T}\lambda
    -
    \frac12c^{\mathsf T}Q^{-1}c.
\end{aligned}
\]

Матрица
\[
    AQ^{-1}A^{\mathsf T}
\]
положительно полуопределена, поэтому \(g\) вогнута. Последнее
постоянное слагаемое не влияет на положение двойственного
максимума.

Если двойственное решение \(\lambda^\star\) существует и выполняется
сильная двойственность, прямое решение восстанавливается по формуле
\[
    x^\star
    =
    -Q^{-1}
    \bigl(
        c+A^{\mathsf T}\lambda^\star
    \bigr).
\]

Если
\[
    m\ll n,
\]
двойственная задача имеет существенно меньшую размерность. Вместо
обмена векторами \(x\in\RR^n\) распределённые агенты могут
согласовывать множители
\[
    \lambda\in\RR^m.
\]

\begin{remark}
Для применения сильной двойственности к задаче \((P_Q)\) достаточно,
например, существования строго допустимой точки
\[
    \bar x
    \quad\text{такой, что}\quad
    A\bar x<b.
\]
Положительная определённость \(Q\) обеспечивает строгую выпуклость
целевой функции и единственность прямого решения, но сама по себе не
заменяет условие регулярности ограничений.
\end{remark}

\section{Связь с машинным обучением}

В методе опорных векторов прямая задача содержит вектор параметров
гиперплоскости и по одному ограничению для каждого объекта обучающей
выборки. Переход к двойственной задаче выражает решение через
множители этих ограничений.

В двойственной форме объекты входят только через скалярные
произведения. Это позволяет заменить скалярное произведение
ядровой функцией и получить нелинейный классификатор без явного
построения пространства признаков. Двойственная форма SVM и
ядровый переход рассматриваются далее.

\begin{lecturesummary}
Двойственный разрыв равен
\[
    p^\star-d^\star\ge0
\]
и может быть строго положительным даже для выпуклой задачи. Условие
Слейтера требует существования строго допустимой точки для выпуклых
неравенств при точном выполнении аффинных равенств. При этом условии
выполняется сильная двойственность:
\[
    p^\star=d^\star.
\]
Оптимальный множитель Лагранжа характеризует чувствительность
оптимального значения к изменению ограничения. Для квадратичной
задачи с \(Q\succ0\) прямая переменная исключается явно:
\[
    x(\lambda)
    =
    -Q^{-1}(c+A^{\mathsf T}\lambda),
\]
что приводит к вогнутой квадратичной двойственной задаче по
неотрицательному вектору \(\lambda\).
\end{lecturesummary}